\section{Conclusions}
\label{sec:conclusion}

Data are a core production factor of the digital economy, but their asset value is difficult to observe. This paper turns the broad question ``what are data worth?'' into a narrow empirical question: when a privacy regime raises the cost of third-party data, how much market value do exposed firms lose? We answer this question using the 2021 ATT and PIPL shock, a text-based TDDI constructed from 2019 annual reports, and a firm-quarter panel of 298 Chinese A-share digital firms from 2018Q1 to 2023Q4.

The evidence supports H1, H2, and H3. The main findings are threefold. First, high-TDDI firms experience a 9.2\% relative decline in log Tobin's Q and a 2.45 percentage-point decline in gross margin after the shock, while ROA and selling-expense ratios do not change significantly. Second, the effect is concentrated in B2B data-service firms, which decline by 14.2\%, while consumer internet firms are not significantly affected. Third, the implied market-value loss for high-exposure firms is about RMB 203.7 billion, with a 95\% confidence interval of RMB 54.7--352.8 billion. The FPDCI split does not support a buffering role for first-party data (H4); instead, first-party data capability and third-party data dependence appear to coexist.

The paper contributes a replicable measure of third-party data dependence, a valuation estimate for a privacy-regime shock, and a data risk exposure that can inform data-asset accounting and privacy policy. The B2B--B2C contrast shows that data assets are not homogeneous: first-party data and third-party data have different substitutability under privacy regulation, and the firms that intermediate external data bear a disproportionate share of the adjustment cost.

Three extensions follow naturally. First, the TDDI can be applied to other markets and other privacy regimes, which would test the external validity of the exposure measure. Second, linking the index to firm-level data flows, app-level tracking, or advertising outcomes would provide a more direct validation of the mechanism. Third, a dynamic model of first-party data investment would clarify whether the observed valuation loss is permanent or is offset by adaptation over time. These extensions would deepen the understanding of how data assets are priced in the digital economy.
